\label{sec:walls}
Work on $\R^n\setminus\{0\}$; everything is conical. The depth-$L$ tiling stratifies input space into relatively open polyhedral cones:
open tiles, open pieces of walls, and so on. These are the \emph{faces}. Each face is convex, hence contractible.
\begin{definition}[Wall foliation]
The wall foliation $\cF_L$ is the module of (formal) vector fields tangent to every wall $Z_{l,a}$, $l\le L$. Its leaves are the faces.
At a point of a face $S$ of codimension $j$ the \emph{transverse model} is the induced formal foliation of a $j$-dimensional slice.
\end{definition}
The radial (Euler) field $x\cdot\nabla$ is tangent to every wall, because walls are cones. So the drift of the Ornstein--Uhlenbeck generator
$\Delta-x\cdot\nabla$ is longitudinal and only the diffusion is transverse. The Euler--Stein identity $\E f=\E[x\cdot\nabla f]=\E[\Delta f]$ for
$1$-homogeneous $f$ records where the longitudinal part meets the transverse one.

\begin{theorem}[Normal and bent crossings]\label{thm:bent}
Let $S$ be a codimension-two face at a generic point of which exactly two walls meet, $Z=Z_{l,a}$ and $Z'=Z_{k,b}$ with $k\le l$. Let
$T=\partial z_{l,a}/\partial h_{k,b}$ be the gated transport coefficient (intermediate gates fixed at the point; $T=0$ if $k=l$).
\begin{enumerate}[leftmargin=*]
\item If $T=0$ the walls are transverse hyperplanes near $S$. The transverse model is the log foliation of $y_1y_2=0$, generated by $y_1\partial_1$
and $y_2\partial_2$; the isotropy Lie algebra is $\R^2$ (abelian). The foliation is Debord near $S$ (its algebroid is the log tangent
bundle), its Nash blowup is trivial, and its Helffer--Nourrigat cone is the whole log cotangent bundle \cite{Louis}. The holonomy groupoid
is the $b$-groupoid of Melrose's calculus \cite{Melrose}.
\item If $T\ne0$, choose transverse coordinates $y_1=z_{k,b}$ and $y_2=$ the linear function equal to $z_{l,a}$ on the side $y_1<0$. Then
$z_{l,a}=y_2+T\relu(y_1)$ near $S$, so $Z=\{y_2=0,\ y_1\le0\}\cup\{y_2+Ty_1=0,\ y_1\ge0\}$ is \emph{bent} along $Z'$. The formal transverse model is
the log foliation of the central three-line arrangement
\[
y_1\,y_2\,(y_2+Ty_1)=0 ,
\]
a free module (Saito) with basis $E=y_1\partial_1+y_2\partial_2$ and $\vartheta_T=-y_1(y_2+Ty_1)\partial_1$. The isotropy Lie algebra is
$\mathrm{span}(E,\vartheta_T)$ with $[E,\vartheta_T]=\vartheta_T$, i.e.\ $\aff(1)$, non-abelian, and the linear isotropy is $\R E$, of dimension $1$
instead of $2$.
\item The integrated isotropy at a bent crossing is the connected $ax+b$ group $\mathrm{Aff}_+(1)$, which Francis identifies with the jet
group $J^2$ \cite{Francis}. The groupoid of $\cF_L$ restricted to $S$ is $S\times S\times\mathrm{Aff}_+(1)$, with $C^*$-algebra $C^*(\mathrm{Aff}_+(1))\otimes\mathcal K$, an
extension of $C_0(\R)$ by $\mathcal K\oplus\mathcal K$. At a normal crossing it is $C_0(\R^2)\otimes\mathcal K$.
\end{enumerate}
\end{theorem}
\begin{proof}
Near a generic point of $S$ only the gate $g_{k,b}$ changes, and crossing $Z'$ toggles $h_{k,b}$ between $0$ and $z_{k,b}$, which changes
$z_{l,a}$ by $T\,z_{k,b}$. This gives the normal form. A formal vector field is tangent to a half-plane through $S$ iff each homogeneous
component is, and a homogeneous polynomial field tangent to a half-line in the slice is tangent to the whole line (polynomial identity on
a ray). So the formal module is the module of fields tangent to the three lines $\{y_1=0\}$, $\{y_2=0\}$, $\{y_2+Ty_1=0\}$; the extensions of the
two half-lines of $Z$ are \emph{phantom} lines that the formal structure cannot distinguish from the wall. Saito's criterion holds:
$\det(E,\vartheta_T)=y_1y_2(y_2+Ty_1)$. $[E,\vartheta_T]=(2-1)\vartheta_T$ because $\vartheta_T$ is homogeneous of degree $2$, and $\vartheta_T\notin I_0\cF$ (otherwise
$\vartheta_T=\ell E$ with $\ell$ linear and the determinant would vanish). A linear field tangent to three distinct lines through $0$ has three
eigenlines, hence is scalar. For $T=0$ two of the lines coincide and the reduced arrangement is $y_1y_2=0$. (3): the free module is the
section module of a Lie algebroid with injective anchor on the regular part, so the holonomy groupoid is the Lie groupoid integrating it
\cite{AS,Louis}; the simply connected group with Lie algebra $\aff(1)$ is $\mathrm{Aff}_+(1)$, which has trivial centre. The structure of
$C^*(\mathrm{Aff}_+(1))$ is classical.
\end{proof}
The check (Verification, item 5) solves the tangency conditions for homogeneous polynomial fields in the slice. The dimensions in degrees
$1,2,3$ are $2,4,6$ at $T=0$ and $1,3,5$ at $T=0.3$, and $[E,\vartheta_T]=\vartheta_T$ holds symbolically.

\paragraph{The bent crossings are the memory.} $T$ is the coefficient with which the kink of neuron $(k,b)$ enters the deep pre-activation:
the coupling at which Price's theorem gives birth to the third cumulant (stage 10, Theorem \emph{birth}). In the simplest case it is literally
the cumulant. For $z_{2,a}=\sum_bW_2(a,b)\relu(w_b\cdot x)$ with orthogonal first-layer rows, the $w_b\cdot x$ are independent, the bending
coefficients are $T_{ab}=W_2(a,b)$, and
\[
\kappa_3(z_{2,a})=c_3\sum_bT_{ab}^3|w_b|^3,\qquad c_3=\kappa_3(\relu Z)=\frac{\pi+2}{(2\pi)^{3/2}}=0.3265 .
\]
So $z_{2,a}$ is Gaussian iff its wall is nowhere bent, and its skewness is the sum of cubed bendings over its bent crossings. With
correlated rows, cross terms appear, weighted by the codegrees (cosines) of the normal crossings among the shallow walls. These are
the walk sums with exact contacts of stage 10, now read as an expansion over the codimension-two strata of the wall foliation: bent
crossings are vertices with a contact, normal crossings are edges. The Gaussian closure is exact for any wall system with only normal
crossings, and every family of jointly Gaussian pre-activations is the pre-activation family of such a system in an augmented Gaussian
space (realize the covariance as $\tilde z_l=M_l\xi$). Its error is therefore carried by the bent crossings.

\paragraph{Smooth versus piecewise linear.} The piecewise-linear homeomorphism $(y_1,y_2)\mapsto(y_1,y_2+T\relu(y_1))$ unbends $Z$ into a
normal crossing. In the PL category every crossing is normal. The bending, the phantom line and the non-abelian isotropy exist only for
the smooth (formal) structure, and the Gaussian measure is a smooth object, not a PL one. This is the geometric content of ``Gaussian
calculus sees the kinks'': non-Gaussianity is the obstruction to smoothing the unbending chart.

\begin{proposition}[No holonomy]\label{prop:noholonomy}
(1) Every face is contractible, so by Fischer--Laurent-Gengoux (\cite{FLG}, Cor.~3.10) the formal wall foliation along any face is the
product of the face with its transverse model: outer holonomy and principal-bundle data are trivial. (2) Every wall $Z_{l,a}$ has the
global defining function $z_{l,a}$; for each $k$ the transverse-order-$k$ structure it induces is Scott's $b^k$ structure with a global
$(k-1)$-jet \cite{Scott}, whose holonomy invariant is trivial (\cite{Francis}, \S1). Hence the noncommutativity of the wall foliation is purely
isotropic: abelian at normal crossings, $\aff(1)$ at bent crossings, and solvable at higher-codimension faces.
\end{proposition}
The contrast with tilings is instructive. In a substitution tiling the boundary relations $R^j_B$ of Julien--Savinien are large: dense and
of tail type. In the network every face is a convex cone and the global defining functions kill all holonomy. Whatever the network has
that is genuinely noncommutative lives at points (isotropy) and in the depth-and-localization dynamics (Sections~\ref{sec:parabolic},
\ref{sec:deformation}), not in the leaf space.
